% Presentation slides. Uses the SAME auto-generated tables, figures and numbers as the
% report (symlinks: tables/, figures/, generated/ -> ../report/...).
% Build with XeLaTeX (Korean via kotex): latexmk -xelatex slides.tex  (or: tectonic slides.tex)
% Speaker notes: add [notes=show] / [notes=only] via \setbeameroption{show notes on second screen=right}
\documentclass[aspectratio=169,11pt]{beamer}
\usepackage{amsmath,amssymb}
\usepackage{kotex}
\setmainhangulfont{UnDotum}[Extension=.ttf]
\setsanshangulfont{UnDotum}[Extension=.ttf]
\usepackage{booktabs,multirow,adjustbox,array}
\usepackage{tikz}
\usetikzlibrary{arrows.meta,positioning,fit,calc}

\usetheme{default}
\usecolortheme{default}
\definecolor{accent}{HTML}{2a78d6}
\definecolor{ink}{HTML}{1f1f1e}
\setbeamercolor{frametitle}{fg=ink}
\setbeamercolor{title}{fg=ink}
\setbeamercolor{structure}{fg=accent}
\setbeamerfont{frametitle}{series=\bfseries}
\setbeamertemplate{navigation symbols}{}
\setbeamertemplate{itemize items}[circle]
\setbeamertemplate{footline}{\hfill\scriptsize\color{gray}\insertframenumber/\inserttotalframenumber\hspace{3mm}\vspace{2mm}}

\newcommand{\TBD}[1]{\textcolor{red}{\textbf{[TBD:\,#1]}}}
\input{generated/numbers}

\title{Training Data Thresholds for\\Cost-Effective Korean Sentiment Analysis}
\subtitle{When do neural models beat TF-IDF --- and is it worth the cost?}
\author{\TBD{author names}}
\institute{Practical Machine Learning, Hanyang University (ERICA)}
\date{\today}

\begin{document}

\begin{frame}
  \titlepage
  \note{Introduce the team and the one-line question: how many labelled reviews do you need before a neural model is worth it?}
\end{frame}

\begin{frame}{Label more reviews, or use a bigger model?}
  \begin{itemize}
    \item A small team has a few hundred labelled Korean reviews and must choose:
          \textbf{label more data} or \textbf{use a more complex model}.
    \item TF-IDF + linear models are cheap; neural sequence models need more data and compute.
  \end{itemize}
  \vspace{2mm}
  \begin{block}{Research question}
    For each dataset, what is the \textbf{minimum training size} (500, 1k, 2k, 4k, 8k) at which
    the best neural model beats the best TF-IDF model by \textbf{$\ge$ \crossoverMargin{} macro-F1
    point}, and is it still competitive in \textbf{training time, latency and memory}?
  \end{block}
  \note{Frame the decision a team actually faces. Simple baselines are known to win at small
  scale (Zhang et al.\ 2015; Joulin et al.\ 2017); we measure where that stops for Korean reviews.}
\end{frame}

\begin{frame}{Three Korean review domains, fixed splits}
  \small
  \textbf{NSMC} movie reviews (\nCleanNSMC{} clean) \quad
  \textbf{Shopping} product reviews, ratings 1--2 / 4--5 (\nCleanShopping) \quad
  \textbf{Steam} game reviews (\nCleanSteam)\\[1mm]
  Fixed 8k pool / 2k val / 2k test per dataset, stratified, no text overlap.
  \vspace{3mm}

  \centering
  \input{tables/table1}
  \note{Three domains with different label sources: movie ratings, shopping stars (3 dropped),
  Steam recommendation flags. All splits are fixed before any modelling.}
\end{frame}

\begin{frame}{Preprocessing for Korean}
  \small
  \begin{columns}[T]
    \begin{column}{0.5\linewidth}
      \textbf{Normalization} (fixed rules, before splitting)
      \begin{itemize}\small
        \item NFC, \emph{not} NFKC (keeps ㅋ ㅠ ㅎ intact)
        \item cap repeats: ㅋㅋㅋㅋㅋ $\to$ ㅋㅋ, !!!!! $\to$ !!
        \item keep digits (``10점 만점에 1점''), emoji, profanity
        \item no spacing correction, no stopwords
      \end{itemize}
    \end{column}
    \begin{column}{0.48\linewidth}
      \textbf{Kiwi morphemes} (\texttt{form/TAG})
      \begin{itemize}\small
        \item 봤는데 $\to$ 보/VV 었/EP 는데/EC
        \item keep particles \& endings: 안/MAG, -지 않-, -네
        \item symbols survive: \symbolShareNSMC{} of NSMC reviews
        \item OOV at $n$=500: \oovSmallNSMC{} $\to$ \oovLargeNSMC{} at 8k
      \end{itemize}
    \end{column}
  \end{columns}
  \vspace{3mm}
  Same tokens for all nine methods; ablation swaps Kiwi for eojeols / characters.
  \note{Explain why agglutination makes eojeol vocabularies sparse, and why sklearn's default token pattern would drop 안/좋/못.}
\end{frame}

\begin{frame}{Nine pipelines, two families}
  \begin{columns}[T]
    \begin{column}{0.48\linewidth}
      \textbf{TF-IDF} (1--2-grams, 10k features)
      \begin{itemize}\small
        \item Logistic regression, linear SVM
        \item Random Forest, XGBoost
      \end{itemize}
    \end{column}
    \begin{column}{0.48\linewidth}
      \textbf{Neural} (embeddings from scratch)
      \begin{itemize}\small
        \item MLP (DAN), BiLSTM, BiGRU
        \item small Transformer, small Mamba
        \item non-embedding params matched: 50k / 200k / 500k
      \end{itemize}
    \end{column}
  \end{columns}
  \vspace{4mm}
  \small Shared Kiwi morphemes; all vocabularies fitted on the training subset only (neural: min.\ frequency 2).
  Same training rules for every size (AdamW, batch 64, early stopping on val macro-F1).
  \note{Four TF-IDF and five neural models; neural backbones are budget-matched so that the
  comparison is about architecture, not size.}
\end{frame}

\begin{frame}{Protocol}
  \centering
  \input{figures/fig1_pipeline}
  \vspace{3mm}
  \small
  \begin{itemize}
    \item 10 predefined configs per method, tuned once at $n$=2k on \textbf{validation}
    \item \nFinalFits{} final fits: 3 datasets $\times$ 9 methods $\times$ 5 sizes $\times$ \nSeeds{} seeds
    \item Paired comparison: same seed $\Rightarrow$ same training reviews for every method
    \item Tokenizer ablation: +\nAblationFits{} fits
  \end{itemize}
  \note{Tuning happens once, on validation, at 2k; configurations are then frozen. Test data is
  never used for any choice.}
\end{frame}

\begin{frame}{Hardware and protocol: parallel for quality, exclusive for cost}
  \small
  \begin{columns}[T]
    \begin{column}{0.5\linewidth}
      \textbf{One machine}
      \begin{itemize}\small
        \item \gpuModel, \cpuModel
        \item fp32; \nThreads{} CPU threads per TF-IDF job and all CPU inference
        \item Neural training on \trainDevice{} (\neuralThreads{} threads); Mamba: \mambaBackend
      \end{itemize}
      \textbf{Two run modes}
      \begin{itemize}\small
        \item \emph{Quality}: \kGpu{} neural jobs share the GPU, \cpuJobs{} TF-IDF jobs alongside
        \item \emph{Cost}: one job at a time, idle machine: all times and memory
        \item Predictions identical in both modes: \detIdentical/\detRuns{} runs
      \end{itemize}
    \end{column}
    \begin{column}{0.48\linewidth}
      \textbf{Compute (busy wall time)}\\[1mm]
      \begin{tabular}{lrr}
        \toprule
        Stage & hours & GPU util \\
        \midrule
        Tuning (exclusive) & \wallTuning & \gpuUtilTuning \\
        Main sweep (parallel) & \wallSweep & \gpuUtilSweep \\
        Ablation (parallel) & \wallAblation & \gpuUtilAblation \\
        Cost pass (exclusive) & \wallCostPass & \gpuUtilCostPass \\
        \midrule
        Total incl.\ pilot & \wallTotal & \\
        \bottomrule
      \end{tabular}
    \end{column}
  \end{columns}
  \note{Running jobs concurrently is fast but distorts timings, so every reported time comes from
  a separate exclusive run. The determinism check shows concurrency never changed a prediction.
  The whole study took \wallTotal{} GPU-hours.}
\end{frame}

\begin{frame}{TF-IDF leads at every size; neural models close the gap}
  \centering
  \includegraphics[width=\linewidth]{figures/fig2_learning_curves.pdf}\\[-1mm]
  {\scriptsize Test macro-F1, mean $\pm$ std over \nSeeds{} seeds (sample std). Dashed = TF-IDF, solid = neural.}
  \note{At 500 reviews the best neural model trails by \ablationGapSmallNSMC{} (NSMC),
  \ablationGapSmallShopping{} (Shopping) and \ablationGapSmallSteam{} (Steam) points. Neural curves
  are steeper, but they never cross. Test OOV drops from \oovSmallNSMC{} to \oovLargeNSMC{} on
  NSMC, one plausible reason for the slow neural start.}
\end{frame}

\begin{frame}{No crossover up to 8,000 reviews in any domain}
  \centering
  \input{tables/table5}
  \vspace{3mm}
  \small
  \begin{itemize}
    \item NSMC: \crossoverNSMC \quad Shopping: \crossoverShopping \quad Steam: \crossoverSteam
    \item Winners chosen on validation; sensitivity (test-selected): \crossoverTestSelNSMC, \crossoverTestSelShopping, \crossoverTestSelSteam
  \end{itemize}
  \note{This answers the research question. The gap at 8k is \gapNSMC, \gapShopping{} and
  \gapSteam{} points, i.e.\ the neural winner is still behind. Steam's neural winner is
  formally Pareto-efficient only because its CPU latency is slightly lower.}
\end{frame}

\begin{frame}{Same millisecond latency, far more training cost}
  \centering
  \includegraphics[width=\linewidth]{figures/fig3_quality_vs_cost.pdf}\\[-1mm]
  {\scriptsize $n$ = 8,000. CPU end-to-end latency per review (batch 1, median); stars = accelerator (\accelDevice) latency; shaded = CPU Pareto front.}
  \note{Per-review CPU latency is about 1 ms for both families; the GPU does not help at batch
  size 1. Training is the real cost: the neural winners need \fitRatioSteam--\fitRatioShopping{}
  the fit time and about \ramRatioNSMC{} the RAM. Random Forest is the latency outlier.}
\end{frame}

\begin{frame}{Kiwi morphemes help both families; the conclusion holds}
  \centering
  \includegraphics[width=0.7\linewidth]{figures/fig5_tokenizers.pdf}\\[-1mm]
  {\scriptsize Best TF-IDF and best neural method per dataset with eojeol / Kiwi / character tokens; mean $\pm$ std over \nSeeds{} seeds.}\\
  \small Gap with each family's best tokenizer at 8k: \ablationGapLargeNSMC{} / \ablationGapLargeShopping{} / \ablationGapLargeSteam{}
  \quad conclusion holds: \ablationHoldsNSMC{} / \ablationHoldsShopping{} / \ablationHoldsSteam
  \note{Character n-grams do not make TF-IDF stronger here; Kiwi is best for both families in
  almost every cell. Shopping at 8k: validation picks character TF-IDF, which is worse on test,
  so the best-tokenizer gap turns positive (\ablationGapLargeShopping), still below
  \crossoverMargin{} point.}
\end{frame}

\begin{frame}{Steam is hardest; on Shopping the families err differently}
  \centering
  \includegraphics[width=\linewidth]{figures/fig4_confusion.pdf}\\[-1mm]
  {\scriptsize Confusion matrices at $n$ = 8,000, summed over \nSeeds{} runs, row-normalised (\%).}
  \note{NSMC: both classes similar. Shopping: TF-IDF misses positives (\fnrTfidfShopping), the
  neural winner mislabels negatives (\fprNeuralShopping), which hints at ensemble potential.
  Steam: about a quarter of each class is wrong for both families; positives slightly harder.}
\end{frame}

\begin{frame}{Conclusion: use TF-IDF + a linear model, and label more data}
  \begin{itemize}
    \item Crossover thresholds: NSMC \crossoverNSMC, Shopping \crossoverShopping, Steam \crossoverSteam.
    \item Best at 8k: \bestMethodNSMC{} (\macroFOneBestNSMC), \bestMethodShopping{} (\macroFOneBestShopping), \bestMethodSteam{} (\macroFOneBestSteam).
    \item Neural winner Pareto-efficient? NSMC \paretoNSMC, Shopping \paretoShopping, Steam \paretoSteam.
    \item No method has saturated: the last doubling still adds \gainLastNSMC{} / \gainLastShopping{} / \gainLastSteam{} points (TF-IDF winner).
    \item \textbf{Recommendation:} up to 8k labelled reviews and no pretrained encoder, deploy a
          linear model on Kiwi TF-IDF features and invest in labels.
  \end{itemize}
  \note{Neural models get closer with more data, and a crossover beyond 8k is plausible but
  untested. Mamba and the Transformer underperform the simple embedding models in this
  small, from-scratch regime.}
\end{frame}

\begin{frame}{Limitations}
  \begin{itemize}
    \item Rating-based labels are imperfect sentiment proxies.
    \item Three domains, binary labels; no pretrained models.
    \item Thresholds resolved only to a factor of two (tested sizes).
    \item Kiwi brings external linguistic knowledge (applied to all methods; see ablation).
    \item Ablation reuses Kiwi-tuned hyperparameters.
    \item One machine (\cpuModel); timings depend on hardware and implementation (Mamba: \mambaBackend).
    \item \nSeeds{} seeds measure variability --- not a significance test.
  \end{itemize}
  \note{The biggest caveat is the lack of pretrained Korean encoders; with them the answer could
  be very different.}
\end{frame}

\begin{frame}{Demo}
  \begin{itemize}
    \item Type a Korean review, pick a dataset: see the normalized text, Kiwi tokens, all
          nine models' predictions, confidences and measured latency.
    \item Run: \texttt{python source/demo/app.py [--device cpu|auto]} (models trained on $n$ = 8,000, seed 13).
  \end{itemize}
  \note{Try a sarcastic review and one with only emoticons (e.g.\ repeated laughing characters).}
\end{frame}

\begin{frame}[plain]
  \centering \Large Thank you --- questions?
\end{frame}

\appendix
\begin{frame}{Backup: cost at $n$ = 8,000}
  \centering
  \scriptsize
  \begin{adjustbox}{max totalheight=0.82\textheight}\input{tables/table4b}\end{adjustbox}
\end{frame}

\begin{frame}{Backup: test macro-F1 across sizes}
  \centering
  \tiny
  \begin{adjustbox}{max totalheight=0.82\textheight}\input{tables/table3}\end{adjustbox}
\end{frame}

\end{document}
